\documentclass{situationist-preprint}

\title{The Case of Insects}
\msid{MND-23-c}
\msfamily{MND-23}
\msdiscipline{Philosophy of Mind}
\mstype{Consequence}
\mskeywords{insects; animal consciousness; sentience; markers; evidential verdicts}

\begin{document}
\maketitle

\begin{abstract}
MND-23-a concludes that the evidence for insect consciousness is inconclusive. We derive what this verdict amounts to and why it is the right one. It is neither a probable affirmative nor a probable negative: some markers are present, others absent or untested, and the positive findings admit explanations without experience that have not been excluded. The verdict is also not robust to the choice of theory, since theories that locate consciousness in midbrain integration favor insects and theories that require cortical or higher-order structures disfavor them. We state what evidence would change it.
\end{abstract}

\section{The verdict and its grounds}

By Definition~3.4 of MND-23-a, a verdict is \emph{inconclusive} when the markers are partial or unconvergent and an explanation without consciousness has not been excluded. Lemma~7.1 applies this to insects. The positive evidence is real: bumblebees show a dopamine-dependent positive bias after unexpected reward \citep{Solvi2016} and roll balls without reward in a way that meets criteria for play \citep{Dona2022}; bees learn by observation \citep{Chittka2022}; the central complex integrates information across modalities \citep{Barron2016}. But wound-directed behavior and analgesic effects are weakly attested, and the evidence differs across orders \citep{Gibbons2022}. The structures are analogous, not homologous, so Principle~P3 of MND-23-a lends no support.

\begin{proposition}\label{p:neither}
The evidence supports neither the verdict that insects are probably conscious nor the verdict that they are probably not.
\end{proposition}

\begin{proof}
An affirmative verdict requires strong convergence (Principle~P2 of MND-23-a); the insect markers are present in some families and missing or untested in others. A negative verdict requires that an explanation without consciousness be the better one; the deflationary explanations of the positive findings, by reinforcement learning and global modulation of state, have been offered but not shown to account for the data better, and the absence of some markers reflects the absence of experiments rather than negative results.
\end{proof}

The verdict is also theory-sensitive in a way the verdicts for mammals and birds are not. Proposition~9.1 of MND-23-a states that those verdicts are robust across theories and that, for insects, the theories diverge. For insects, a theory on which consciousness depends on midbrain-type integration \citep{Merker2007,Barron2016} favors ascription, whereas theories requiring cortical or higher-order structures disfavor it. A theory-light method cannot choose between them, and so cannot go beyond an inconclusive verdict.

\begin{objection}[Precaution settles it]
Given the stakes, insects should be counted as conscious.
\end{objection}

\begin{reply}
Precaution bears on what to do under uncertainty, not on what the evidence shows \citep{Birch2024}. It may justify protective measures for insects whose sentience is a realistic possibility; it does not convert an inconclusive verdict into a probable one.
\end{reply}

\section{What would change it}

The verdict would move toward ascription if wound-directed behavior and analgesic-reversible avoidance were found across orders under controlled conditions, and toward denial if the positive findings were shown to be reproduced by systems known to lack experience with the same profile of flexibility. Until then, the right answer is that we do not know.

\begin{thebibliography}{9}\small
\bibitem[Barron and Klein(2016)]{Barron2016} Barron, A.~B. and Klein, C. (2016). What insects can tell us about the origins of consciousness. \emph{Proceedings of the National Academy of Sciences}, 113(18), 4900--4908.
\bibitem[Birch(2024)]{Birch2024} Birch, J. (2024). \emph{The Edge of Sentience: Risk and Precaution in Humans, Other Animals, and AI}. Oxford University Press.
\bibitem[Chittka(2022)]{Chittka2022} Chittka, L. (2022). \emph{The Mind of a Bee}. Princeton University Press.
\bibitem[Dona et~al.(2022)]{Dona2022} Dona, H.~S.~G., Solvi, C., Kowalewska, A., M\"akel\"a, K., MaBouDi, H. and Chittka, L. (2022). Do bumble bees play? \emph{Animal Behaviour}, 194, 239--251.
\bibitem[Gibbons et~al.(2022)]{Gibbons2022} Gibbons, M., Crump, A., Barrett, M., Sarlak, S., Birch, J. and Chittka, L. (2022). Can insects feel pain? A review of the neural and behavioural evidence. \emph{Advances in Insect Physiology}, 63, 155--229.
\bibitem[Merker(2007)]{Merker2007} Merker, B. (2007). Consciousness without a cerebral cortex: A challenge for neuroscience and medicine. \emph{Behavioral and Brain Sciences}, 30(1), 63--81.
\bibitem[Solvi et~al.(2016)]{Solvi2016} Solvi, C., Baciadonna, L. and Chittka, L. (2016). Unexpected rewards induce dopamine-dependent positive emotion-like state changes in bumblebees. \emph{Science}, 353(6307), 1529--1531.
\end{thebibliography}

\end{document}
